\section{Criteri di selezione dei casi di test}

\subsection{Criteri adottati}
% TODO: quali criteri si usano per scegliere i dati di test (classi di equivalenza, valori limite, copertura) e perché; cenno ad affidabilità e validità del criterio.

\subsection{Partizioni di input}
% TODO: per le funzionalità principali, tabella delle classi di equivalenza (valide e non valide) e dei valori limite.

\begin{center}
\begin{tblr}{width=\linewidth, colspec={X[1,l]X[1.5,l]X[1.5,l]X[1.5,l]}, cells={font=\small},
             row{1}={bg=arancio}, hline{2-Z}={wd=0.5pt, fg=arancio}}
Input & Classi valide & Classi non valide & Valori limite \\
 & & & \\
\end{tblr}
\end{center}

\subsection{Criteri di copertura e di uscita}
% TODO: soglie di copertura del codice da raggiungere e condizioni per considerare concluso il testing.
